\chapter{Administrer la plateforme}

\section{Accueil Super Admin}

L'accueil Super Admin présente l'état global de la plateforme : organisations,
utilisateurs, demandes en attente, activité commerciale et événements
nécessitant une action.

\manualscreenshot{superadmin-home.png}
  {Vue d'ensemble de la plateforme}
  {fig:superadmin-home}

\section{Organisations}

La page \emph{Organizations} regroupe les organisations, leur statut, leur
capacité, leur abonnement et leurs administrateurs.

\manualscreenshot{superadmin-organizations.png}
  {Gestion globale des organisations}
  {fig:superadmin-organizations}

\begin{procedurebox}{Contrôler une organisation}
  \begin{enumerate}[leftmargin=*]
    \item Rechercher l'organisation par son nom ou sa référence.
    \item Vérifier son statut commercial, son plan et ses quotas.
    \item Consulter l'administrateur principal et l'historique récent.
    \item Effectuer uniquement l'action nécessaire : mise à jour commerciale,
    suspension, restauration ou assistance.
    \item Vérifier la présence de la trace correspondante dans l'audit.
  \end{enumerate}
\end{procedurebox}

\section{Demandes de démonstration}

Une demande de démonstration est destinée à une future organisation. Le
demandeur n'obtient pas immédiatement un espace administrateur.

\manualscreenshot{superadmin-demo-requests.png}
  {Traitement des demandes de démonstration}
  {fig:superadmin-demo}

\subsection{Workflow}

\begin{enumerate}[leftmargin=*]
  \item Le visiteur transmet ses coordonnées, son organisation et plusieurs
  objectifs éventuels.
  \item La plateforme accuse réception sans exposer de lien interne de revue.
  \item Le Super Admin examine la demande et choisit de la rejeter ou de
  poursuivre.
  \item Lors de l'approbation, il configure l'organisation, la durée d'essai et
  le compte administrateur initial.
  \item La plateforme crée l'organisation isolée, sa souscription d'essai et
  l'invitation de l'administrateur.
  \item Le demandeur active son compte, définit son mot de passe et complète
  éventuellement le logo de son organisation.
\end{enumerate}

Le lien d'activation de l'administrateur est valable 48 heures par défaut et
peut être réglé globalement entre 12 et 168 heures.

\attention{Renvoyer une invitation n'est pertinent que pour une invitation
encore attendue ou arrivée à expiration. Une invitation révoquée doit être
explicitement recréée selon le workflow prévu.}

\section{Plans et décisions commerciales}

Le Super Admin définit les offres, capacités, cycles de facturation et règles de
renouvellement. Une demande de changement formulée par un administrateur doit
être examinée avant de modifier la souscription.

\begin{procedurebox}{Restaurer une organisation suspendue}
  \begin{enumerate}[leftmargin=*]
    \item Vérifier le motif de suspension et les factures ouvertes.
    \item Enregistrer la décision commerciale ou le règlement simulé/confirmé.
    \item Affecter un plan valide et ses dates d'effet.
    \item Restaurer l'organisation.
    \item Vérifier que les administrateurs sont notifiés et que les quotas sont
    cohérents.
  \end{enumerate}
\end{procedurebox}

\section{Paramètres globaux}

Les paramètres globaux contrôlent notamment l'inscription publique, les
demandes de démonstration, la durée des invitations, l'essai, la période de
grâce, les quotas, la rétention de l'audit et l'adresse de support.

\manualscreenshotcompact{superadmin-settings.png}
  {Paramètres globaux et garde-fous d'environnement}
  {fig:superadmin-settings}

\begin{table}[H]
  \centering
  \caption{Exemples de paramètres globaux}
  \label{tab:platform-settings}
  \begin{tabularx}{\textwidth}{|L{4.8cm}|L{3.2cm}|Y|}
    \hline
    \rowcolor{ctmint}
    \textbf{Paramètre} & \textbf{Plage} & \textbf{Impact} \\ \hline
    Inscription client publique & Activé / désactivé & Contrôle l'accès à la
    page d'inscription locale. \\ \hline
    Demandes de démonstration & Activé / désactivé & Contrôle le formulaire
    public de demande. \\ \hline
    Invitation employé & 12 à 336 heures & Durée de validité du lien
    d'activation. \\ \hline
    Essai organisation & 7 à 90 jours & Durée proposée lors du provisionnement.
    \\ \hline
    Grâce commerciale & 1 à 30 jours & Délai avant suspension opérationnelle.
    \\ \hline
    Rétention de l'audit & 30 à 730 jours & Conservation de l'historique global.
    \\ \hline
  \end{tabularx}
\end{table}

\section{Intégrations}

La page \emph{Integrations} indique l'état de Google OAuth, du service d'envoi
transactionnel, de la passerelle OCPP et de l'adaptateur de paiement. Les
secrets restent dans l'environnement de déploiement et ne sont jamais renvoyés
au navigateur.

\begin{table}[H]
  \centering
  \caption{Intégrations de la version de référence}
  \label{tab:integrations}
  \begin{tabularx}{\textwidth}{|L{3.5cm}|L{4.2cm}|Y|}
    \hline
    \rowcolor{ctmint}
    \textbf{Domaine} & \textbf{Technologie} & \textbf{Usage} \\ \hline
    Authentification & Google OAuth 2.0 & Connexion sociale des comptes clients.
    \\ \hline
    Courriel & Resend & Vérification, invitation, récupération et notifications.
    \\ \hline
    Recharge & OCPP 1.6 JSON & Connexion des bornes et du simulateur SAP. \\ \hline
    Paiement & Adaptateur WireMock & Simulation des réponses et webhooks du
    prestataire. \\ \hline
  \end{tabularx}
\end{table}

\section{Audit}

Le journal d'audit permet de rechercher l'auteur, l'action, l'entité, la date et
le contexte d'une opération sensible. Il sert au diagnostic et à la
responsabilisation, mais ne remplace ni une sauvegarde ni les journaux
techniques du serveur.

\securitynote{Un Super Admin ne doit jamais saisir un secret dans une note, un
rapport ou un champ audité. Les secrets Google, Resend, OCPP et paiement restent
dans les variables d'environnement protégées.}
